\documentclass{article}%
\usepackage[T1]{fontenc}%
\usepackage[utf8]{inputenc}%
\usepackage{lmodern}%
\usepackage{textcomp}%
\usepackage{lastpage}%
\usepackage{expex}%
\usepackage{geometry}%
\geometry{margin=2cm}%
%
\lingset{everygla=\it}%
\title{Análisis Lingüístico: Gramatica{-}Normativa{-}Kaqchikel.pdf}%
\author{Páginas: 171{-}171}%
%
\begin{document}%
\normalsize%
\maketitle%
\section{Page 171}%
\label{sec:Page171}%

%
(plural) del objeto y sujeto. La palabra principal incluye primeramente el tiempo/aspecto %
o modo (excepto en el perfectivo que es la última), seguido por el juego absolutivo, %
juego ergativo, la raíz o base en calidad de obligatorios y sufijos de categoría en casos %
específicos (si el verbo es derivado). Como elementos opcionales incluye el movimiento, %
su posición es siempre antes del juego ergativo. Ejemplos: %
\subsection{Marca de tiempo/aspecto/modo:}%
\label{subsec:Marcadetiempo/aspecto/modo}%

%
\ex
  	abla Xb'entz'eta' na kan ri ak'wal.
  	ablz X- 0- b'e- n- tz'et- a' na kan ri ak'wal
  	able \textbf{COM-} \textbf{B3S-} \textbf{MOV-} \textbf{Als-} \textbf{ver-} \textbf{SC} \textbf{PAR} \textbf{DIR} \textbf{DET} \textbf{niño}
  	ablf [Xb'entz'eta']_{SVT}
  \xe[\textit{Todavía fui a ver al niño.}]
%
\subsection{Marca el objeto en el verbo:}%
\label{subsec:Marcaelobjetoenelverbo}%

%
\ex
  	abla ...kan jeb'ël ok xerutz'ët ri ch'uta'q xtani' rija'...
  	ablz kan jeb'ël ok x- e- ru- tz'ët ri ch'uta'q xtan- i' rija'
  	able \textbf{PAR} \textbf{hermosa} \textbf{DIR} \textbf{COM-} \textbf{B3p-} \textbf{A3s-} \textbf{ver} \textbf{DET} \textbf{DIM} \textbf{Niña-} \textbf{PL} \textbf{él}
  	ablf [xerutz'ët]_{SVT}
  \xe[\textit{...vio muy hermosas a las niñas...}]
%
\subsection{Marca el sujeto en el verbo:}%
\label{subsec:Marcaelsujetoenelverbo}%

%
\ex
  	abla Xkilöq' jun uqaj ri ka'i' achi'a'.
  	ablz X- 0- ki- löq jun uqaj ri ka'i' achi'- a'.
  	able \textbf{COM-} \textbf{B3s-} \textbf{A3p-} \textbf{comprar} \textbf{IND} \textbf{corte} \textbf{DET} \textbf{dos} \textbf{hombre-} \textbf{PL}
  	ablf [Xkilöq']_{SVT}
  \xe[\textit{Los dos hombres compraron un corte.}]
%
\ex
  	abla Xusipaj ri xkoya' ri xtän.
  	ablz X- 0- u- sip- a- j ri xkoya' ri xtän.
  	able \textbf{COM-} \textbf{B3s-} \textbf{A3s-} \textbf{regalar-} \textbf{BV-} \textbf{SC} \textbf{DET} \textbf{tomate} \textbf{DET} \textbf{señorita}
  	ablf [Xusipaj]_{SVT}
  \xe[\textit{La señorita regaló el tomate.}]
%
\subsection{Marca movimiento en el verbo:}%
\label{subsec:Marcamovimientoenelverbo}%

%
\ex
  	abla ...nb'esamäj qatata', nb'eqato'....
  	ablf [nb'eqato']_{SVT}
  \xe[\textit{...nuestro papá va a trabajar, vamos a ayudarlo....}]
%
\end{document}